\section{Hypotheses in Full}
\label{sup:hyp}
Each hypothesis states what would falsify it, and several state where the gate is not expected to
help. Thresholds match the pre-registered configuration files.
\begin{description}
  \item[H1] \emph{(residuals first; go/no-go).} Fit $\sigma(\cdot)$, $\ell$, $\nu_{\mathrm{nug}}$, the tails, the
        shared distance-bias coefficient $\beta$, the range of the scene temperature $\gamma$ and
        the per-view heterogeneity $\omega$ of the aligned Depth Pro residuals on UseGeo and ETH3D
        development scenes; characterise natural regional failures, separately for nadir views from
        survey altitude and for near-ground views. If $\omega\gtrsim0.05$, the region scan leads
        and the energy test keeps its $\omega$-calibrated threshold. \emph{Go} if $\ge10\%$ of views carry a failure covering $\ge2\%$ of pixels at
        $>3\sigma$ and keeping those views lowers the fused-surface F-score by $\ge2$ points.
        \emph{No-go} (pivot to an analysis paper) below 5\% of views or below 1 point.
  \item[H2] \emph{(repetition).} Cross-view correlation of aligned errors at the same reference
        surface point: $\ge0.5$ on textureless regions, $\le0.3$ on specular, $\le0.2$ on textured, and
        decreasing with the separation of the two frames along the strip. \emph{Falsified} if the
        textureless value is below 0.5 or not at least 0.2 above the specular one. The decay with
        separation sets the block size $B$.
  \item[H3] \emph{(detection).} At 5\% false alarm, on specular-region failures, the gate beats the
        best cheap or robust baseline by $\ge15$ points of power. \emph{Falsified} if the gain is
        below 15 points or its interval includes zero. No gain is predicted on textureless failures.
  \item[H4] \emph{(reconstruction).} With the gate, AbsRel and F-score improve by $\ge10\%$ relative
        over no gate and the best robust-splatting baseline on failure scenes of UseGeo and ETH3D, and
        on MatrixCity with injected faults,
        and change by $\le2\%$ on clean scenes. \emph{Falsified} if the paired interval includes zero
        on failure scenes, or clean scenes degrade by more than 2\%.
  \item[H5] \emph{(calibration).} With independent per-surfel blocks, the raw mean $T_j/m_j$ on
        clean views exceeds 1.2 (overconfidence). With the leave-view-out scene temperature and the
        $\omega$-calibrated energy threshold, the per-view false-alarm rate on test scenes is within
        $[0.5\alpha,1.5\alpha]$ without per-scene retuning; deflation by $n_{\mathrm{eff}}/n$ and
        jackknife inflation are ablations. \emph{Falsified} if it leaves that band.
  \item[H6] \emph{(ablations).} Joint surface $-10$ points; keeping tested-view-only surfels $-20$
        points on large region failures; unwhitened $-5$ points; rendered blocks at least as calibrated
        and powerful as the FisherRF-style and Bayes' Rays-style uncertainties; dropping the scene
        temperature or the shared distance bias raises the clean false-alarm rate; a fold-wise opacity
        warm-up changes the false-alarm rate and the power by less than 2 points, so freezing opacity is
        safe. The outer loop detects at least 10 points more faulty views than a single pass when two or
        more faulty views share a region, with the clean false-alarm rate still in its band; dropping
        all rejected views per round excludes more clean views; readmission changes the results by
        less than 2 points. Leave-block-out with $B$ chosen from the sweep $B\in\{1,3,5,9\}$ detects at least
        10 points more faulty views than $B=1$ when neighbouring frames repeat the error, at a cost in
        untestable pixels that is reported; starting from perturbed poses changes the detection power
        by less than 5 points. \emph{Falsified} if any ablation is within its margin, if the fold-wise
        warm-up changes the results by more, if the loop does not beat the single pass by 10 points,
        or if the chosen block does not beat $B=1$ by 10 points.
  \item[H7] \emph{(pose perturbation check).} After the gate has run, the poses of chosen views (isolated
        views, a block, and the flagged view) are perturbed by known rotations
        ($0.05$--$2^{\circ}$) and translations (1--50 ground sampling distances). (a)~Before refinement, with
        their own depth excluded from the tested pose, the gate flags at least 80\% of the views
        perturbed by at least $0.5^{\circ}$ or 10 ground sampling distances. (b)~Spline, pose-graph and
        render-and-compare refinement removes at least 80\% of the perturbation up to $1^{\circ}$ and 25
        ground sampling distances. Recovery error is reported against perturbation size, including where
        it fails. \emph{Falsified} if either rate is below 80\% in its range.
\end{description}
H1 is a go/no-go and runs first; H2 predicts a blind spot rather than a success; H3 and H4 are the
claims that would justify the method; H5 tests the stated approximation; H6 tests the design choices; H7 is a controlled check of the
pose refinement and of the gate on pose errors.

\paragraph{Breakdown scenarios} Reported as honest negatives: errors repeated in all views
(textureless surfaces, H2), including across a whole block; low overlap, where pixels are untestable; faulty views that supply most of a
region; masking and swamping across rounds of the outer loop, and loops that reach
the round cap; a misspecified noise shape or a per-view heterogeneity above the calibrated $\omega$, which the
scene temperature cannot fix; and faults or misfits beyond the linearised fusion.

\section{Limitations and Future Work in Full}
\label{sup:limits}
\subsection{Limitations}
(1)~Independent per-surfel blocks ignore cross-surfel correlations, the correlation of depth error
within a view and the shared pose error of a view, so the rendered variance is overconfident; we do not
fix this here. The threshold is calibrated on clean data, H5 measures the gap, and deflation by
$n_{\mathrm{eff}}/n$ or jackknife inflation are ablations only. (2)~The scene temperature corrects the
noise level, not its shape or tails, and not a view whose nominal noise is genuinely larger than the
scene's; such a view keeps an elevated false-alarm rate and is indistinguishable from a fault by design,
so the guarantee is a scene-average false-alarm rate. (3)~If the depth model repeats the same error
across views, as H2 tests for textureless surfaces, the other views agree with the fault and the gate
cannot see it; leave-block-out removes the neighbours that share it, but not errors repeated across
whole strips, and it weakens the cavity, raises the detection floor and leaves more pixels untestable. (4)~Pixels that no other view sees are untestable, so low overlap lowers availability.
(5)~If faulty views are a majority in a region, the leave-view-out surface follows the fault and the
outer loop locks onto the wrong consensus; masking among a minority of faulty views and swamping of
clean views are reduced, not excluded, by dropping one view per round, and a dropped clean view is not
readmitted by default. (6)~The leave-out is a Newton step at the MAP, exact only for the linear model,
and opacity is not re-estimated without the tested view; within a round, a surfel pruned because of a
not-yet-caught view cannot be restored, although the next round re-densifies once that view is caught.
The Holm schedule over rounds is exact only for fixed p-values, so the loop's level rests on
clean-run calibration. (7)~Whole-view affine errors are absorbed by the alignment and are neither
detected nor harmful; the depth bias is modelled only as a shared distance term. (8)~Colour may refine the
poses but never the map geometry, so appearance cues are not used for detection. (9)~Pose faults are
not a target of the gate, only of the controlled check in H7; poses are refined, not estimated from scratch, and the spline assumes a smooth,
continuously recorded flight. (10)~The noise shape is calibrated on a few static scenes for one default
model, and nadir aerial imagery lies outside the depth model's main training distribution, which H1
measures rather than assumes away.

\subsection{Future Work}
Planned, none started: modelling cross-surfel correlation by low-rank or banded blocks; testing poses
as well as depth; a per-view bias field of the depth model, which is the subject of a separate paper; using the rendered uncertainty for view selection and for reconstruction from fewer
views; online operation; and depth from physical sensors. Estimating the light sources and relighting the
appearance block, for example on aerial data under varying illumination \cite{data_beyondlight}, could later
let colour help without letting it shape the geometry.
